\textbf{Keywords:} Range, Antenna Size, Modulation, OFDM, GFSK,
pi/4-qpsk, 8-psk, 16 QAM, Bluetooth LE, LP RX, LNA, Mixer, AAF, ADC, BB

Radio's are all around us. In our phone, on our wrist, in our house,
there is Bluetooth, WiFi, Zigbee, LTE, GPS and many more.

A radio is a device that receives and transmits light encoded with
information. The frequency of the light depends on the standard. How the
information is encoded onto the light depends on the standard.

Assume that we did not know any standards, what would we do if we wanted
to make the best radio IC for gaming mice?

There are a few key concepts we would have to know before we decide on a
radio type: Data Rate, Carrier Frequency and range, and the power
supply.

\section{Data Rate}\label{data-rate}\index{data rate@Data rate}

\subsection{Data}\label{data}\index{data@Data}

A mouse reports on the relative X and Y displacement of the mouse as a
function of time. A mouse has buttons. There can be many mice in a room,
as such, they must have an address , so PCs can tell them apart.

A mouse must be low-power. As such, the radio cannot be on all the time.
The radio must start up and be ready to receive quickly.

We don't know how far away from the PC the mice might be, as such, we
don't know the dB loss in the communication channel. As a result, the
radio needs to have a high dynamic range, from weak signals to strong
signals. In order for the radio to adjust the gain of the receiver we
should include a pre-amble, a known sequence, for example 01010101, such
that the radio can adjust the gain, and also, recover the symbol timing.

All in all, the packets we send from the mouse may need to have the
following bits.

\begin{longtable}[]{@{}lll@{}}
\toprule\noalign{}
What & Bits & Why \\
\midrule\noalign{}
\endhead
\bottomrule\noalign{}
\endlastfoot
X displacement & 8 & \\
Y displacement & 8 & \\
CRC & 4 & Bit errors \\
Buttons & 16 & One-hot coding. Most mice have buttons \\
Preamble & 8 & Synchronization \\
Address & 32 & Unique identifier \\
Total & 76 & \\
\end{longtable}

\subsection{Rate}\label{rate}\index{rate@Rate}

Gamers are crazy for speed, they care about milliseconds. So our mice
needs to be able to send and receive data quite often.

Assume 1 ms update rate

\subsection{Data Rate}\label{data-rate-1}\index{data rate@Data rate}

To compute the data rate, let's do a back of the envelope estimate of
the data, and the rate.

\begin{Shaded}
\begin{Highlighting}[]
\ExtensionTok{Application}\NormalTok{ Data Rate }\OperatorTok{\textgreater{}}\NormalTok{ 76 bits/ms = 76 kbps}

\ExtensionTok{Assume}\NormalTok{ 30 \% packet loss}

\ExtensionTok{Raw}\NormalTok{ Data Rate }\OperatorTok{\textgreater{}}\NormalTok{ 228 kbps}

\ExtensionTok{Multiply}\NormalTok{ by 3.14 }\OperatorTok{\textgreater{}}\NormalTok{ 716 kbps}

\ExtensionTok{Round}\NormalTok{ to nearest nice number = 1Mbps}
\end{Highlighting}
\end{Shaded}

The above statements are a exact copy of what happens in industry when
we start design of something. We make an educated guess and multiply by
a number. More optimistic people would multiply with \(e\).

\section{Carrier Frequency \& Range}\label{carrier-frequency-range}

\subsection{ISM (industrial, scientific and medical)
bands}\label{ism-industrial-scientific-and-medical-bands}

There are rules and regulations that prevent us from transmitting and
receiving at any frequency we want. We need to pick one of the ISM
bands, or we need to get a license from governments around the world.

For the ISM bands, there are regions, as seen below.


{
\centering
\aicfig{media/International_Telecommunication_Union_regions_with_dividing_lines.pdf}
}

\emph{Figure 1: ITU regions that set the ISM band allocations: region 1
(yellow), region 2 (blue), region 3 (pink). Image: Maximilian
Doerrbecker (Chumwa), CC BY-SA 2.5, via Wikimedia Commons}

\begin{itemize}
\tightlist
\item
  Yellow: Region 1
\item
  Blue: Region 2
\item
  Pink: Region 3
\end{itemize}

Below is a table of the available frequencies, but how should we pick
which one to use? There are at least two criteria that should be
investigated. Antenna and Range.

\begin{longtable}[]{@{}llll@{}}
\toprule\noalign{}
Flow & Fhigh & Bandwidth & Description \\
\midrule\noalign{}
\endhead
\bottomrule\noalign{}
\endlastfoot
40.66 MHz & 40.7 MHz & 40 kHz & Worldwide \\
433.05 MHz & 434.79 MHz & 1.74 MHz & Region 1 \\
902 MHz & 928 MHz & 26 MHz & Region 2 \\
2.4 GHz & 2.5 GHz & 100 MHz & Worldwide \\
5.725 GHz & 5.875 GHz & 150 MHz & Worldwide \\
24 GHz & 24.25 GHz & 250 MHz & Worldwide \\
61 GHz & 61.5 GHz & 500 MHz & Subject to local acceptance \\
\end{longtable}

\subsection{Antenna}\label{antenna}\index{antenna@Antenna}

For a mouse we want to hold in our hand, there is a size limit to the
antenna. There are many types of antenna, but

assume wavelength/4 is an OK antenna size (wavelength =
lightspeed/frequency)

The below table shows the ISM band and the size of a quarter wavelength
antenna. Any frequency above 2.4 GHz may be OK from a size perspective.

\begin{longtable}[]{@{}lrrr@{}}
\toprule\noalign{}
ISM band & \(\lambda/4\) & Unit & OK/NOK \\
\midrule\noalign{}
\endhead
\bottomrule\noalign{}
\endlastfoot
40.68 MHz & 1.8 & m & :x: \\
433.92 MHz & 17 & cm & :x: \\
915 MHz & 8.2 & cm & \\
2450 MHz & 3.06 & cm & :white\_check\_mark: \\
5800 MHz & 1.29 & cm & :white\_check\_mark: \\
24.125 GHz & 3.1 & mm & :white\_check\_mark: \\
61.25 GHz & 1.2 & mm & :white\_check\_mark: \\
\end{longtable}

\subsection{Range (Friis)}\label{range-friis}\index{range (friis)@Range (friis)}

One of the worst questions a radio designer can get is ``What is the
range of your radio?'', especially if the people asking are those that
don't understand physics, or the real world. The answer to the question
is incredibly complicated, as it depends on exactly what is between two
devices talking.

If we assume, however, that there is only free space, and no real
reflections from anywhere, then we can make an estimate of the range.

Assume no antenna gain, power density p at distance D is

\[ p = \frac{P_{TX}}{4 \pi D^2}\]

Assume receiver antenna has no gain, then the effective aperture is

\[ A_e = \frac{\lambda^2}{4 \pi}\]

Power received is then

\[P_{RX} = \frac{P_{TX}}{D^2} \left[\frac{\lambda}{4 \pi}\right]^2\]

Or in terms of distance

\[ D = 10^\frac{P_{TX} - P_{RX} + 20 log_{10}\left(\frac{c}{4 \pi f}\right)}{20} \]

\subsection{Range (Free space)}\label{range-free-space}\index{range (free space)@Range (free space)}

If we take the ideal equation above, and use some realistic numbers for
TX and RX power, we can estimate a range.

Assume TX = 0 dBm, assume RX sensitivity is -80 dBm

\begin{longtable}[]{@{}lcrr@{}}
\toprule\noalign{}
Freq & \textbf{\(20 log_{10}\left(c/4 \pi f\right)\)} {[}dB{]} & D
{[}m{]} & OK/NOK \\
\midrule\noalign{}
\endhead
\bottomrule\noalign{}
\endlastfoot
915 MHz & -31.7 & 260.9 & :white\_check\_mark: \\
\textbf{2.45 GHz} & \textbf{-40.2} & \textbf{97.4} &
:white\_check\_mark: \\
5.80 GHz & -47.7 & 41.2 & :white\_check\_mark: \\
24.12 GHz & -60.1 & 9.9 & :x: \\
61.25 GHz & -68.2 & 3.9 & :x: \\
160 GHz & -76.52 & 1.5 & :x: \\
\end{longtable}

\subsection{Range (Real world)}\label{range-real-world}\index{range (real world)@Range (real world)}

In the real world, however, the

path loss factor, \[ n \in [1.6,6]\],
\[ D = 10^\frac{P_{TX} - P_{RX} + 20 log_{10}\left(\frac{c}{4 \pi f}\right)}{n
\times 10} \]

So the real world range of a radio can vary more than an order of
magnitude. Still, 2.4 GHz seems like a good choice for a mouse.

\begin{longtable}[]{@{}
  >{\raggedright\arraybackslash}p{(\linewidth - 8\tabcolsep) * \real{0.1818}}
  >{\centering\arraybackslash}p{(\linewidth - 8\tabcolsep) * \real{0.2727}}
  >{\raggedleft\arraybackslash}p{(\linewidth - 8\tabcolsep) * \real{0.1818}}
  >{\raggedleft\arraybackslash}p{(\linewidth - 8\tabcolsep) * \real{0.1818}}
  >{\raggedleft\arraybackslash}p{(\linewidth - 8\tabcolsep) * \real{0.1818}}@{}}
\toprule\noalign{}
\begin{minipage}[b]{\linewidth}\raggedright
Freq
\end{minipage} & \begin{minipage}[b]{\linewidth}\centering
\textbf{\[20 log_{10}\left(c/4 \pi f\right)\]} {[}dB{]}
\end{minipage} & \begin{minipage}[b]{\linewidth}\raggedleft
D@n=2 {[}m{]}
\end{minipage} & \begin{minipage}[b]{\linewidth}\raggedleft
D@n=6 {[}m{]}
\end{minipage} & \begin{minipage}[b]{\linewidth}\raggedleft
OK/NOK
\end{minipage} \\
\midrule\noalign{}
\endhead
\bottomrule\noalign{}
\endlastfoot
\textbf{2.45 GHz} & \textbf{-40.2} & \textbf{97.4} & \textbf{4.6} &
:white\_check\_mark: \\
5.80 GHz & -47.7 & 41.2 & 3.45 & :white\_check\_mark: \\
24.12 GHz & -60.1 & 9.9 & 2.1 & :x: \\
\end{longtable}

\section{Power supply}\label{power-supply}\index{power supply@Power supply}

We could have a wired mouse for power, but that's boring. Why would we
want a wired mouse to have wireless communication? It must be powered by
a battery, but what type of battery?

There exists a bible of batteries, Linden's Handbook of Batteries. It's
worth a read if you want to dive deeper into chemistry and properties of
primary (non-chargeable) and secondary (chargeable) cells.

\subsection{Battery}\label{battery}\index{battery@Battery}

Mouse is maybe AA, 3000 mAh

\begin{longtable}[]{@{}llrr@{}}
\toprule\noalign{}
Cell & Chemistry & Voltage (V) & Capacity (Ah) \\
\midrule\noalign{}
\endhead
\bottomrule\noalign{}
\endlastfoot
AA & LiFeS2 & 1.0 - 1.8 & 3 \\
2xAA & LiFeS2 & 2.0 - 3.6 & 3 \\
AA & Zn/Alk/MnO2 & 0.8 - 1.6 & 3 \\
2xAA & Zn/Alk/MnO2 & 1.6 - 3.2 & 3 \\
\end{longtable}

\section{Decisions}\label{decisions}

Now we know that we need a 1 Mbps radio at 2.4 GHz that runs off a 1.0 V
- 1.8 V or 2.0 V - 3.6 V supply.

Next we need to decide what modulation scheme we want for our light. How
should we encode the bits onto the 2.4 GHz carrier wave?

\subsection{Modulation}\label{modulation}\index{modulation@Modulation}

Any modulation can be described by the function below.

\[ A_m(t) \times \cos\left( 2 \pi \int_0^t f_{carrier}(\tau)d\tau + \phi_{m}(t)\right)\]

The integral matters as soon as the carrier frequency is one of the
things being modulated, which for GFSK it is. Writing \(2\pi f(t)t\)
instead is a tempting shorthand and it is wrong: differentiate it and
the instantaneous frequency comes out as \(f(t) + t\,f'(t)\), so the
error grows without bound as \(t\) does. Phase is the integral of
frequency, always. For a fixed carrier the integral collapses to the
familiar \(2\pi f_c t\) and no harm is done, which is why the shorthand
survives.

The amplitude of the carrier can be modulated, or the phase of the
carrier.

People have been creative over the last 50 years in terms of encoding
bits onto carriers. Below is a small excerpt of some common schemes.

\begin{longtable}[]{@{}
  >{\raggedright\arraybackslash}p{(\linewidth - 6\tabcolsep) * \real{0.2500}}
  >{\raggedright\arraybackslash}p{(\linewidth - 6\tabcolsep) * \real{0.2500}}
  >{\raggedright\arraybackslash}p{(\linewidth - 6\tabcolsep) * \real{0.2500}}
  >{\raggedright\arraybackslash}p{(\linewidth - 6\tabcolsep) * \real{0.2500}}@{}}
\toprule\noalign{}
\begin{minipage}[b]{\linewidth}\raggedright
Scheme
\end{minipage} & \begin{minipage}[b]{\linewidth}\raggedright
Acronym
\end{minipage} & \begin{minipage}[b]{\linewidth}\raggedright
Pro
\end{minipage} & \begin{minipage}[b]{\linewidth}\raggedright
Con
\end{minipage} \\
\midrule\noalign{}
\endhead
\bottomrule\noalign{}
\endlastfoot
Binary phase shift keying & BPSK & Simple & Not constant envelope \\
Quadrature phase-shift keying & QPSK & 2bits/symbol & Not constant
envelope \\
Offset QPSK & OQPSK & 2bits/symbol & Constant envelope with half-sine
pulse shaping \\
Gaussian Frequency Shift Keying & GFSK & 1 bit/symbol & Constant
envelope \\
Quadrature amplitude modulation & QAM & \textgreater{} 10 bits/symbol &
Really non-constant envelope \\
\end{longtable}

\subsection{BPSK}\label{bpsk}\index{bpsk@BPSK}

In binary phase shift keying the 1 and 0 is encoded in the phase change.
Change the phase 180 degrees and we've transitioned from a 0 to a 1. Do
another 180 degrees and we're back to where we were.

It's common to show modulation schemes in a constellation diagram with
the real axis and the complex axis. For the real signal we send, the
phase and amplitude are usually both real quantities.

I say usually, because in quantum mechanics, and the time evolution of a
particle, the amplitude of the wave function is actually a complex
variable. As such, nature is actually complex at the most fundamental
level.

But for now, let's keep it real in the real world.

Still, the maths is much more elegant in the complex plane.

The equation for the unit circle is \(y = e^{i( \omega t + \phi)}\)
where \(\phi\) is the phase, and \(\omega\) is the angular frequency.

Imagine we spin a bike wheel around at a constant frequency (constant
\(\omega\)), on the bike wheel there is a red dot. If you keep your eyes
open all the time, then the red dot would go round and round. But
imagine that you only opened your eyes every second for a brief moment
to see where the dot was. Sometimes it could be on the right side,
sometimes on the left side. If our ``eye opening rate'', or your sample
rate, matched how fast the ``wheel rotator'' changed the location of the
dot, then you could receive information.

Now imagine you have a strobe light matched to the ``normal'' carrier
frequency. If one rotation of the wheel matched the frequency of the
strobe light, then the red dot would stay in exactly the same place. If
the wheel rotation was slightly faster, then the red dot would move one
way around the circle at every strobe. If the wheel rotation was
slightly slower, the red dot would move the other way around the circle.

That's exactly how we can change the position in the constellation. We
increase the carrier frequency for a bit to rotate 180 degrees, and we
can decrease the frequency to go back 180 degrees. In this example the
dot would move around the unit circle, and the amplitude of the carrier
can stay constant.


{
\centering
\aicfig{media/l7_bpsk_real_tikz.pdf}
}

\emph{Figure 2: BPSK constellation: the two symbols sit on the real
axis, 180 degrees apart}

There is another way to change phase 180 degrees, and that's simply to
swap the phase in the transmitter circuit. Imagine as below we have a
local oscillator driving pseudo differential common source stages with
switches on top. If we flip the switches we can change the phase 180
degrees pretty fast.

A challenge is, however, that the amplitude will change. In general,
constant envelope (don't change amplitude) modulation is less bandwidth
efficient (slower) than schemes that change both phase and amplitude.


{
\centering
\aicfig{media/l7_bpsk_circuit_tikz.pdf}
}

\emph{Figure 3: BPSK transmitter: a local oscillator drives a
pseudo-differential common source pair, and the \(b_0\) switches swap
the output phase 180 degrees into the antenna balun}

Standards like Zigbee used offset quadrature phase shift keying, with a
constellation as shown below. With 4 points we can send 2 bits per
symbol.


{
\centering
\aicfig{media/l7_qpsk_tikz.pdf}
}

\emph{Figure 4: QPSK constellation: four symbols at \(\pm 1 \pm j\), or
\(\sqrt{2}e^{\pm j\pi/4}\), so 2 bits per symbol}

In ZigBee, or 802.15.4 as the standard is called, the phase changes is
actually done with a constant envelope.

The nice thing about constant envelope is that the radio transmitter can
be simple. We don't need to change the amplitude. If we have a PLL as a
local oscillator, where we can change the phase (or frequency), then we
only need a power amplifier before the antenna.


{
\centering
\aicfig{media/l7_const_env_tikz.pdf}
}

\emph{Figure 5: Constant envelope transmitter: the phase is modulated in
the local oscillator, and a power amplifier drives the antenna}

For phase and amplitude modulation, or complex transmitters, we need a
way to change the amplitude and phase. What a shocker. There are two
ways to do that. A polar architecture where phase change is done in the
PLL, and amplitude in the power amplifier.


{
\centering
\aicfig{media/l7_polar_tikz.pdf}
}

\emph{Figure 6: Polar transmitter: phase \(\phi\) is modulated in the
local oscillator and amplitude \(A\) in the power amplifier}

Or a Cartesian architecture where we make the in-phase component, and
quadrature-phase components in digital, then use two digital to analog
converters, and a set of complex mixers to encode onto the carrier. The
power amplifier would not need to change the amplitude, but it does need
to be linear.


{
\centering
\aicfig{media/l8_cartesian_tikz.pdf}
}

\emph{Figure 7: Cartesian transmitter. Two converters produce \(I\) and
\(Q\), two mixers multiply them by copies of the local oscillator ninety
degrees apart, and the sum drives a linear power amplifier. The
\(90^\circ\) block is the whole difference between this and two copies
of the same branch: without it both mixers would be multiplying by the
same carrier, and the sum would carry no phase information.}

We can continue to add constellation points around the unit circle.
Below we can see 8-PSK, where we can send 3-bits per symbol. Assuming we
could shift position between the constellation points at a fixed rate,
i.e 1 mega symbols per second. With 1-bit per symbol we'd get 1 Mbps.
With 3-bits per symbol we'd get 3 Mbps.

We could add 16 points, 32 points and so on to the unit circle, however,
there is always noise in the transmitter, which will create a cloud
around each constellation point, and it's harder and harder to
distinguish the points from each other.


{
\centering
\aicfig{media/l8_8psk_tikz.pdf}
}

\emph{Figure 8: 8-PSK constellation: eight points on the unit circle, so
3 bits per symbol}

Bluetooth Classic uses \(\pi/4\)-DQPSK and 8DPSK.

DPSK means differential phase shift keying. Think about DPSK like this.
In the QPSK diagram above the symbols (00,01,10,11) are determined by
the constellation point \(1 + j\), \(1-j\) and so on. What would happen
if the constellation rotated slowly? Would \(1+j\) turn into \(1-j\) at
some point? That might screw up our decoding if the received
constellation point was at \(1 + 0j\), we would not know what it was.

If we encoded the symbols as a change in phase instead (differential),
then it would not matter if the constellation rotated slowly. A change
from \(1+j\) to \(1-j\) would still be 90 degrees.

Why would the constellation rotate you ask? Imagine the transmitter
transmits at 2 400 000 000 Hz. How does our receiver generate the same
frequency? We need a reference and a PLL. The crystal-oscillator
reference has a variation of +-50 ppm, so \(2.4e9 \times 50/1e6 = 120\)
kHz.

Assume our receiver local oscillator was at 2 400 120 000 Hz. The
transmitter sends 2 400 000 000 Hz + modulation. At the receiver we
multiply with our local oscillator, and if you remember your math,
multiplication of two sine creates a sum and a difference between the
two frequencies. As such, the low frequency part (the difference between
the frequencies) would be 120 kHz + modulation. As a result, our
constellation would rotate 120 000 times per second. Assuming a symbol
rate of 1MS/s our constellation would rotate roughly 1/10 of the way
each symbol.

In DPSK the rotation is not that important. In PSK we have to measure
the carrier offset, and continuously de-rotate the constellation.

Most radios will de-rotate somewhat based on the preamble, for example
in Bluetooth Low Energy there is an initial 10101010 sequence that we
can use to estimate the offset between TX and RX carriers, or the
frequency offset.

The \(\pi/4\) part of \(\pi/4-DQPSK\) just means we actively rotate the
constellation 45 degrees every symbol, as a consequence, the amplitude
never goes through the origin. In the transmitter circuit, it's
difficult to turn the carrier off, so we try to avoid the zero point in
the constellation.


{
\centering
\aicfig{media/EDR.png}
}

\emph{Figure 9: Bluetooth Core Specification, radio specification:
Enhanced Data Rate uses \(\pi/4\)-DQPSK at 2 Mb/s and 8DPSK at 3 Mb/s}

I don't think 16PSK is that common, at 4-bits per symbol it's common to
switch to Quadrature Amplitude Modulation (QAM), as shown below. The
goal of QAM is to maximize the distance between each symbol. The
challenge with QAM is the amplitude modulation. The modulation scheme is
sensitive to variations in the transmitter amplitude. As such, more
complex circuits than 8PSK could be necessary.

If you wanted to research ``new fancy modulation schemes'' I'd think
about \href{https://en.wikipedia.org/wiki/Sphere_packing}{Sphere
packing}.


{
\centering
\aicfig{media/l8_16qam_tikz.pdf}
}

\emph{Figure 10: 16-QAM constellation: a 4 by 4 grid of points in phase
and amplitude, so 4 bits per symbol}

\subsection{Single carrier, or multi
carrier?}\label{single-carrier-or-multi-carrier}

Assume we wanted to send 1 Gbps over the air. We could choose a
bandwidth of about 1 GHz with 1 bit per symbol, or a bandwidth of 100
MHz if we sent 1024 QAM at 100 MS/s. Both cases would look like the
figure below.

In both cases we get problems with the physical communication channel,
the change in phase and amplitude affect what is received. For a 1 GHz
bandwidth at 2.4 GHz carrier we'd have problems with the phase. At 1024
QAM we'd have problems with the amplitude.


{
\centering
\aicfig{media/l10_single_carrier_tikz.pdf}
}

\emph{Figure 11: Single carrier link: amplitude \(A_m(t)\) and phase
\(\phi_m(t)\) are modulated onto I and Q, transmitted, and de-modulated
after the receiver}

Back in 1966
\href{https://en.wikipedia.org/wiki/Orthogonal_frequency-division_multiplexing\#:~:text=OFDM\%20is\%20a\%20frequency\%2Ddivision,is\%20divided\%20into\%20multiple\%20streams.}{Orthogonal
frequency division multiplexing} was introduced to deal with the
communication channel. In OFDM we modulate a number of sub-carriers in
the frequency space with our wanted modulation scheme (BPSK, PSK, QAM),
then do an inverse fourier transform to get the time domain signal, mix
on to the carrier, and transmit. At the receiver we take an FFT and do
demodulation in the frequency space. See example in figure below.

The name ``multiple carriers'' is a bit misleading. Although there are
multiple carriers on the left and right side of the figure, there is
normally still just one carrier in the TX/RX.


{
\centering
\aicfig{media/l10_multiple_carrier_tikz.pdf}
}

\emph{Figure 12: OFDM link: the sub-carriers are modulated in frequency
space, an IFFT makes the time domain I and Q for the transmitter, and an
FFT at the receiver recovers the sub-carriers}

There are more details in OFDM than the simple statement above, but the
details are just to fix challenges, such as ``How do I recover the
symbol timing? How do I correct for frequency offset? How do I ensure
that my time domain signal terminates correctly for every FFT chunk''

The genius with OFDM is that we can pick a few of the sub-carriers to be
pilot tones that carry no new information. If we knew exactly what was
sent in phase and amplitude, then we could measure the phase and
amplitude change due to the physical communication channel, and we could
correct the frequency space before we tried to de-modulate.

It's possible to do the same with single carrier modulation also.
Imagine we made a 128-QAM modulation on a single carrier. As long as we
constructed the time domain signal correctly (cyclic prefix to make the
FFT work nicely, some preamble to measure the communication channel,
then we could take an FFT at the receiver, correct the phase and
amplitude, do an IFFT and demodulate the time-domain signal as normal.

In radio design there are so many choices it's easy to get lost.

\subsection{Use a Software Defined
Radio}\label{use-a-software-defined-radio}

For our mouse, what radio scheme should we choose? One common instance
of ``how to make a choice'' in industry is ``Delay the choice as long as
possible so you're sure the choice is right''.

Maybe the best would be to use a software defined radio receiver?
Something like the picture below, an antenna, low noise amplifier, and a
analog-to-digital converter. That way we could support any transmitter.
Fantastic idea, right?


{
\centering
\aicfig{media/lg_lna_adc_tikz.pdf}
}

\emph{Figure 13: Software defined radio receiver: antenna, low noise
amplifier and analog-to-digital converter, nothing else}

Well, lets check if it's a good idea. We know we'll use 2.4 GHz, so we
need about 2.5 GHz bandwidth, at least. We know we want good range, so
maybe 100 dB dynamic range. In analog to digital converter design there
are figure of merits, so we can actually compute a rough power
consumption for such an ADC.

ADC FOM \[ = \frac{P}{2 BW 2^n}\]

State of the art FOM \[\approx 5 \text{ fJ/step}\]

\[ BW = 2.5\text{ GHz}\]

\[ DR = 100\text{ dB} \Rightarrow \text{Bits} = (100-1.76)/6.02 \approx 16\text{ bit} \]

\[ P = 5\text{ fJ/step} \times 5 \text{ GHz} \times 2^{16} = 1.6\text{ W}\]

At 1.6 W our mouse would only last for 2 hours. That's too short. It
will never be a low power idea to convert the full 2.5 GHz bandwidth to
digital, we need some bandwidth selectivity in the receive chain.

\section{Bluetooth}\label{bluetooth}\index{bluetooth@Bluetooth}

\href{https://www.bluetooth.com/specifications/specs/core-specification-5-4/}{Bluetooth}
was made to be a ``simple'' standard and was introduced in 1998. The
standard has continued to develop, with Low Energy introduced in 2010.
The latest planned changes can be seen at
\href{https://www.bluetooth.com/specifications/specifications-in-development/}{Specifications
in Development}.

\subsection{Bluetooth Basic Rate/Extended Data
rate}\label{bluetooth-basic-rateextended-data-rate}

\begin{itemize}
\tightlist
\item
  2.400 GHz to 2.4835 GHz
\item
  1 MHz channel spacing
\item
  78 Channels
\item
  Up to 20 dBm
\item
  Minimum -70 dBm sensitivity (1 Mbps)
\item
  1 MHz GFSK (1 Mbps), pi/4-DQPSK (2 Mbps), 8DPSK (3 Mbps)
\end{itemize}

You'll find BR/EDR in most audio equipment, cars and legacy devices. For
new devices (also audio), there is now a transition to Bluetooth Low
Energy.

\subsection{Bluetooth Low Energy}\label{bluetooth-low-energy}\index{bluetooth low energy@Bluetooth low energy}

\begin{itemize}
\tightlist
\item
  2.400 GHz to 2.480 GHz
\item
  2 MHz channel spacing
\item
  40 Channels (3 primary advertising channels)
\item
  Up to 20 dBm
\item
  Minimum -70 dBm sensitivity (1 Mbps)
\item
  1 MHz GFSK (1 Mbps, 500 kbps, 125 kbps), 2 MHz GFSK (2 Mbps)
\end{itemize}

Below are the Bluetooth LE channels. The green are the advertiser
channels, the blue are the data channels, and the yellow is the WiFi
channels.

The advertiser channels have been intentionally placed where there is
space between the WiFi channels to decrease the probability of
collisions.


{
\centering
\aicfig{media/map.PNG}
}

\emph{Figure 14: Bluetooth LE channel map: advertising channels (green)
at 2402, 2426 and 2480 MHz, data channels (blue), and the WiFi channels
(yellow)}

Any Bluetooth LE peripheral will advertise its presence, it will wake up
once in a while (every few hundred milliseconds, to seconds) and
transmit a short ``I'm here'' packet. After transmitting it will wait a
bit in receive to see if anyone responds.

A Bluetooth LE central will camp in receive on a advertiser channel and
look for these short messages from peripherals. If one is observed, the
Central may choose to respond.

Take any spectrum analyzer anywhere, and you'll see traffic on 2402,
2426, and 2480 MHz.


{
\centering
\aicfig{media/ble_connection.jpg}
}

\emph{Figure 15: Advertising: the peripheral transmits advertisements
once per advertisement interval while the central scans, until the
central initiates a connection. Source: Nordic Semiconductor DevZone}

In a connection a central and peripheral (the master/slave names below
have been removed from the spec, that was a fun update to a 3500 page
document) will have agreed on an interval to talk. Every ``connection
interval'' they will transmit and receive data. The connection interval
is tunable from 7.5 ms to seconds.

Bluetooth LE is the perfect standard for wireless mice.


{
\centering
\aicfig{media/ble_connection.png}
}

\emph{Figure 16: In a connection, central and peripheral exchange
transmit and receive packets once every connection interval. Source:
Nordic Semiconductor DevZone}

further information
\href{https://devzone.nordicsemi.com/guides/nrf-connect-sdk-guides/b/software/posts/building-a-ble-application-on-ncs-comparing-and-contrasting-to-softdevice-based-ble-applications}{Building
a Bluetooth application on nRF Connect SDK}

\href{https://www.bluetooth.com/specifications/specifications-in-development/}{Bluetooth
Specifications in Development}

\section{Algorithm to design state-of-the-art LE
radio}\label{algorithm-to-design-state-of-the-art-le-radio}

\begin{itemize}
\tightlist
\item
  Find most recent digest from International Solid State Circuit
  Conference (ISSCC)
\item
  Find Bluetooth low energy papers
\item
  Pick the best blocks from each paper
\end{itemize}

A typical Bluetooth radio may look something like the picture below.
There would be a single antenna for both RX and Tx. There will be some
way to combine the transmit and receive path in a match, or balun.

The receive chain would have a LNA, mixer, anti-alias filter and
analog-to-digital converters. It's likely that the receive path would be
complex (in-phase and quadrature phase) after mixer.

There would be a local oscillator (all-digital phase-locked-loop) to
provide the frequency to the mixers and transmit path, which could be
either polar or Cartesian.


{
\centering
\aicfig{media/l10_lprxarch_tikz.pdf}
}

\emph{Figure 17: Typical Bluetooth radio: antenna and match, LNA, mixer,
I and Q anti-alias filters and ADCs, an all-digital PLL, and the
transmit path}

In the typical radio we'll need the blocks below. I've added a column
for how many people I would want if I was to lead development of a new
radio.

\begin{longtable}[]{@{}
  >{\raggedright\arraybackslash}p{(\linewidth - 6\tabcolsep) * \real{0.1959}}
  >{\raggedright\arraybackslash}p{(\linewidth - 6\tabcolsep) * \real{0.4021}}
  >{\raggedright\arraybackslash}p{(\linewidth - 6\tabcolsep) * \real{0.1546}}
  >{\raggedright\arraybackslash}p{(\linewidth - 6\tabcolsep) * \real{0.2474}}@{}}
\toprule\noalign{}
\begin{minipage}[b]{\linewidth}\raggedright
Blocks
\end{minipage} & \begin{minipage}[b]{\linewidth}\raggedright
Key parameter
\end{minipage} & \begin{minipage}[b]{\linewidth}\raggedright
Architecture
\end{minipage} & \begin{minipage}[b]{\linewidth}\raggedright
Complexity (nr people)
\end{minipage} \\
\midrule\noalign{}
\endhead
\bottomrule\noalign{}
\endlastfoot
Antenna & Gain, impedance & lambda/4 & \textless1 \\
RF match & loss, input impedance & PI-match & \textless1 \\
Low noise amp & NF, current, linearity & LNTA & 1 \\
Mixer & NF, current, linearity & Passive & 1 \\
Anti-alias filter & NF, current, linearity & Active-RC & 1 \\
ADC & Sample rate, dynamic range, linearity & NS-SAR & 1 - 2 \\
PLL & Phase noise, current & AD-PLL & 2-3 \\
Baseband & Eb/N0, gate count, current. & SystemVerilog & \textgreater{}
10 \\
\end{longtable}

\subsection{LNTA}\label{lnta}\index{lnta@LNTA}

The first thing that must happen in the radio is to amplify the noise as
early as possible. Any circuit has inherent noise, be it thermal-,
flicker-, burst-, or shot-noise. The earlier we can amplify the input
noise, the less contribution there will be from the radio circuits.

The challenges in the low noise amplifier is to provide the right gain.
If there is a strong input signal, then reduce the gain. If there is a
low input signal, then increase the gain.

One way to implement variable gain is to reconfigure the LNA. For an
example, see

30.5 A 0.5V BLE Transceiver with a 1.9mW RX Achieving -96.4dBm
Sensitivity and 4.1dB Adjacent Channel Rejection at 1MHz Offset in 22nm
FDSOI \citeproc{ref-tamura20}{{[}1{]}}

A typical Low Noise Transconductance Amplifier is seen below. It's a
combination of both a common source, and a common gate amplifier. The
current in the NMOS and PMOS is controlled by Vgp and Vgn. Keep in mind
that at RF frequencies the signals are weak, so it's easy to provide the
DC for the LNA with a resistor to a diode connected PMOS or NMOS.

In a LNA the input impedance must be matched to what is required by the
antenna/match in order to have maximum power transfer, that's the role
of the inductors/capacitors.


{
\centering
\aicfig{media/l10_lna_tikz.pdf}
}

\emph{Figure 18: Low noise transconductance amplifier: complementary
common source PMOS and NMOS, AC coupled from the antenna match and
biased through resistors by \(V_{gp}\) and \(V_{gn}\)}

\subsection{MIXER}\label{mixer}\index{mixer@MIXER}

In the mixer we multiply the input signal with our local oscillator.
Most often a complex mixer is used. There is nothing complex about
complex signal processing, just read

Complex signal processing is not complex
\citeproc{ref-martin04}{{[}2{]}}

In order to reduce power, it's most common with a passive mixer as shown
below. A passive mixer is just MOS that we turn on and off with 25\%
duty-cycle. See example in

A 370uW 5.5dB-NF BLE/BT5.0/IEEE 802.15.4-Compliant Receiver with
\textgreater63dB Adjacent Channel Rejection at \textgreater2 Channels
Offset in 22nm FDSOI \citeproc{ref-thijssen20}{{[}3{]}}


{
\centering
\aicfig{media/l10_mix_tikz.pdf}
}

\emph{Figure 19: Passive complex mixer: four MOS switches driven by 25\%
duty-cycle clocks \(I_1\), \(I_2\), \(Q_1\) and \(Q_2\) split the LNA
current into the I and Q outputs. Each gate is AC coupled to its clock
and biased to \(V_n\) through a resistor, so the switch sees a
rail-to-rail drive while its DC operating point is set independently.
Note in the timing diagram that the four phases abut: four quarters fill
the period exactly, so between them the switches carry the whole of the
LNA current and none of it is thrown away.}

To generate the quadrature and in-phase clock signals, which must be 90
degrees phase offset, it's common to generate twice the frequency in the
local oscillator (4.8 GHz), and then divide down to 4 2.4 GHz clock
signals.

If the LO is the same as the carrier, then the modulation signal will be
at DC, often called direct conversion.

The challenge at DC is that there is flicker noise, offset, and burst
noise. The modulation type, however, can impact whether low frequency
noise is an issue. In OFDM we can choose to skip the sub-carriers around
0 Hz, and direct conversion works well. An advantage with direct
conversion is that there is no ``image frequency'' and we can use the
full complex bandwidth.

For FSK and direct conversion the low frequency noise can cause issues,
as such, it's common to offset the LO from the transmitted signal, for
example 4 MHz offset. The low frequency noise problem disappears,
however, we now have a challenge with the image frequency (-4 MHz) that
must be rejected, and we need an increased bandwidth.

There is no ``one correct choice'', there are trade-offs that both ways.
KISS (Keep It Simple Stupid) is one of my guiding principles when
working on radio architecture.

These days most de-modulation happens in digital, and we need to convert
the analog signal to digital, but first AAF.

\subsection{AAF}\label{aaf}\index{aaf@AAF}

The anti alias filter rejects frequencies that can fold into the band of
interest due to sampling. A simple active-RC filters is often good
enough.

We often need gain in the AAF, as the LNA does not have sufficient gain
for the weakest signals. -100 dBm in 50 ohm is 2.2 \(\mu\)V RMS, while
the input range of an ADC may be 1 V. Assume we place the lowest input
signal at 0.1 V, so we need a voltage gain of
\(20\log(0.1/2.2\times 10^{-6}) \approx 93\) dB in the receiver.


{
\centering
\aicfig{media/l4_activebiquad_tikz.pdf}
}

\emph{Figure 20: General purpose Active-RC biquad, used here as the
anti-alias filter}

\subsection{ADC}\label{adc}\index{adc@ADC}

Aaah, ADCs, an IP close to my heart. I did my Ph.d and Post-Doc on ADCs,
and the Ph.D students I've co-supervised have worked on ADCs.

At NTNU there have been multiple students through the years that have
made world-class ADCs, and there's still students at NTNU working on
state-of-the-art ADCs.

These days, a good option is a SAR, or a Noise-Shaped SAR.

If I were to pick, I'd make something like A 68 dB SNDR Compiled
Noise-Shaping SAR ADC With On-Chip CDAC Calibration
\citeproc{ref-garvik19}{{[}4{]}} as shown in the figure below.


{
\centering
\aicfig{media/l6_harald_arch.png}
}

\emph{Figure 21: Architecture of the noise-shaping SAR ADC: capacitive
DAC with multiplexers, loop filter H(z), integrating comparator, SAR
logic, calibration logic and code correction}


{
\centering
\aicfig{media/l6_fig_harald_circuit.png}
}

\emph{Figure 22: The switched-capacitor loop filter with two OTAs (the
first one chopped), the clock phases relative to the SAR activity, and
the resulting NTF with -27.8 dB in-band suppression}

Or if I did not need high resolution, I'd choose my trusty A Compiled
9-bit 20-MS/s 3.5-fJ/conv.step SAR ADC in 28-nm FDSOI for Bluetooth Low
Energy Receivers \citeproc{ref-wulff17}{{[}5{]}}.

The main selling point of that ADC was that it's compiled from a
\href{https://github.com/wulffern/sun_sar9b_sky130nm/blob/main/cic/ip.json}{JSON}
file, a
\href{https://github.com/wulffern/sun_sar9b_sky130nm/blob/main/cic/ip.spi}{SPICE}
file and a
\href{https://github.com/wulffern/sun_sar9b_sky130nm/blob/main/cic/sky130.tech}{technology}
file into a DRC/LVS clean layout.

I also included a few circuit improvements. The bottom plate of the SAR
capacitor is in the clock loop for the comparator (DN0, DP1 below), as
such, the delay of the comparator automatically adjusts with capacitance
corner, so it's more robust over corners


{
\centering
\aicfig{media/fig_sar_logic.pdf}
}

\emph{Figure 23: SAR ADC schematic: (a) capacitor array with self-timed
SAR logic chain and comparator, (b) enable flip-flop, (c) bottom-plate
switching of the CDAC, (d) comparator clock generation}

The compiled nature also made it possible to quickly change the
transistor technology. Below is a picture with 180 nm FDSOI transistors
on the left, and 28 nm FDSOI transistors on the right.

I detest doing anything twice, so I love the fact that I never have to
re-draw that ADC again. I just fix the technology file (and maybe some
tweaks to the other files), and I have a completed ADC.


{
\centering
\aicfig{media/l06_fig_toplevel.pdf}
}

\emph{Figure 24: Layout of the two compiled SAR ADCs with comparator,
logic, CDAC and switch: (a) 180 nm IO-transistor version, (b)
core-transistor version}

\subsection{AD-PLL}\label{ad-pll}\index{ad-pll@AD-PLL}

The phase locked loop is the heart of the radio, and it's probably the
most difficult part to make. Depends a bit on technology, but these
days, All Digital PLLs are cool. Start by reading Razavi's PLL book.

You can spend your life on PLLs.


{
\centering
\aicfig{media/l08_pll_2mod_tikz.pdf}
}

\emph{Figure 25: Two-point modulation: the modulation is applied to the
oscillator and the opposite signal to the sigma-delta feedback divider,
so the loop does not see it}

AD-PLL with Bang-Bang phase detector for steady-state


{
\centering
\aicfig{media/pll_master_arch_28feb2020.pdf}
}

\emph{Figure 26: All-digital PLL with a bang-bang phase detector for
steady-state: phase error logic, digital loop filter, DCO calibration
engine with frequency offset estimator, and steady-state detect}

\subsection{Baseband}\label{baseband}\index{baseband@Baseband}

Once the signal has been converted to digital, then the de-modulation,
and signal fixing start. That's for another course, but there are
interesting challenges.

\begin{longtable}[]{@{}
  >{\raggedright\arraybackslash}p{(\linewidth - 2\tabcolsep) * \real{0.5000}}
  >{\raggedright\arraybackslash}p{(\linewidth - 2\tabcolsep) * \real{0.5000}}@{}}
\toprule\noalign{}
\begin{minipage}[b]{\linewidth}\raggedright
Baseband block
\end{minipage} & \begin{minipage}[b]{\linewidth}\raggedright
Why
\end{minipage} \\
\midrule\noalign{}
\endhead
\bottomrule\noalign{}
\endlastfoot
Mixer? & If we're using low intermediate frequency to avoid DC offset
problems and flicker noise \\
Channel filters? & If the AAF is insufficient for adjacent channel \\
Power detection & To be able to control the gain of the radio \\
Phase extraction & Assuming we're using FSK \\
Timing recovery & Figure out when to slice the symbol \\
Bit detection & single slice, multi-bit slice, correlators etc \\
Address detection & Is the packet for us? \\
Header detection & What does the packet contain \\
CRC & Does the packet have bit errors \\
Payload de-crypt & Most links are encrypted by AES \\
Memory access & Payload need to be stored until CPU can do something \\
\end{longtable}

\section{What do we really want, in the
end?}\label{what-do-we-really-want-in-the-end}

The receiver part can be summed up in one equation for the sensitivity.
The noise in a certain bandwidth. The Noise Figure of the analog
receiver. The Energy per bit over Noise of the de-modulator.

\[P_{RX_{sens}} = -174 \text{ dBm} + 10 log_{10}(R_b)  + NF + E_b/N_0\]

Term by term: \(-174\) dBm is the thermal noise in one hertz at room
temperature, \(R_b\) is the bit rate, which sets how much bandwidth that
noise is collected over, \(NF\) is what the receiver's own noise adds,
and \(E_b/N_0\) is what the demodulator needs to hit its error rate.
Note that \(R_b\) here is the data rate; earlier in this chapter \(DR\)
meant dynamic range, which is a different quantity entirely.

The useful move is to run it backwards. The nRF5340 datasheet quotes
\(-97.5\) dBm sensitivity at 1 Mbps, and \(10log_{10}(10^6) = 60\) dB,
so

\[ P_{RX_{sens}} + 174 - 60 =  NF + E_b/N_0 = 16.5 \text{ dB}\]

and that 16.5 dB is the entire budget shared between the analog front
end and the demodulator. GFSK needs something like 12 dB of \(E_b/N_0\),
which leaves only a few decibels of noise figure for everything in front
of it. That is the number the rest of this chapter is really about.


{
\centering
\aicfig{media/nrf53_rx.png}
}

\emph{Figure 27: nRF5340 radio specification: -97.5 dBm sensitivity at 1
Mbps Bluetooth LE, 2.6 mA in receive and 3.2 mA in transmit. Source:
Nordic Semiconductor, nRF5340 Product Specification}

In the block diagram of the device the radio might be a small box, and
the person using the radio might not realize how complex the radio
actually is.

I hope you understand now that it's actually complicated.


{
\centering
\aicfig{media/nrf53.png}
}

\emph{Figure 28: nRF5340 block diagram, where the entire radio is the
single RADIO block (circled) in the network core}

\section{Summary}\label{summary}

The one-page version of this chapter:

\begin{itemize}
\tightlist
\item
  Start from the link budget: Friis in free space, a rather worse
  exponent indoors, and the antenna wants its fraction of a wavelength
\item
  The ISM bands set the playing field; 2.4 GHz trades antenna size
  against propagation and company
\item
  Energy per bit is the real currency: modulation choice, data rate and
  duty cycle set the average current, and the battery sets the lifetime
\item
  Single-carrier modulation keeps the PA efficient (constant envelope);
  multi-carrier buys spectral efficiency at the cost of backoff
\item
  The receive chain is LNA, mixer and filter: the LNA sets the noise
  figure, the mixer moves the band, and everything after runs at a
  friendlier frequency
\item
  A software-defined radio on the bench teaches more about radios than
  any equation in this chapter
\end{itemize}

\section{Would you like to know
more?}\label{would-you-like-to-know-more}

A 0.5V BLE Transceiver with a 1.9mW RX Achieving -96.4dBm Sensitivity
and 4.1dB Adjacent Channel Rejection at 1MHz Offset in 22nm FDSOI
\citeproc{ref-tamura20}{{[}1{]}}, M. Tamura, Sony Semiconductor
Solutions, Atsugi, Japan, 30.5, ISSCC 2020

A 370uW 5.5dB-NF BLE/BT5.0/IEEE 802.15.4-Compliant Receiver with
\textgreater63dB Adjacent Channel Rejection at \textgreater2 Channels
Offset in 22nm FDSOI \citeproc{ref-thijssen20}{{[}3{]}}, B. J. Thijssen,
University of Twente, Enschede, The Netherlands

A 68 dB SNDR Compiled Noise-Shaping SAR ADC With On-Chip CDAC
Calibration \citeproc{ref-garvik19}{{[}4{]}}, H. Garvik, C. Wulff, T.
Ytterdal

A Compiled 9-bit 20-MS/s 3.5-fJ/conv.step SAR ADC in 28-nm FDSOI for
Bluetooth Low Energy Receivers \citeproc{ref-wulff17}{{[}5{]}}, C.
Wulff, T. Ytterdal

Cole Nielsen, \url{https://github.com/nielscol/thesis_presentations}

``Python Framework for Design and Simulation of Integer-N ADPLLs'', Cole
Nielsen,
\url{https://github.com/nielscol/tfe4580-report/blob/master/report.pdf}

Design of CMOS Phase-Locked Loops \citeproc{ref-razavi20}{{[}6{]}},
Behzad Razavi, University of California, Los Angeles

\phantomsection\label{refs}
\begin{CSLReferences}{0}{0}
\bibitem[\citeproctext]{ref-tamura20}
\CSLLeftMargin{{[}1{]} }%
\CSLRightInline{M. Tamura \emph{et al.}, {``30.5 a 0.5V BLE transceiver
with a 1.9mW RX achieving \(-\)96.4dBm sensitivity and 4.1dB adjacent
channel rejection at 1MHz offset in 22nm FDSOI,''} in \emph{2020 IEEE
international solid-state circuits conference - (ISSCC)}, 2020, pp.
468--470, doi:
\href{https://doi.org/10.1109/ISSCC19947.2020.9063021}{10.1109/ISSCC19947.2020.9063021}.
}

\bibitem[\citeproctext]{ref-martin04}
\CSLLeftMargin{{[}2{]} }%
\CSLRightInline{K. Martin, {``Complex signal processing is not
complex,''} \emph{IEEE Transactions on Circuits and Systems I: Regular
Papers}, vol. 51, no. 9, pp. 1823--1836, 2004, doi:
\href{https://doi.org/10.1109/TCSI.2004.834522}{10.1109/TCSI.2004.834522}.
}

\bibitem[\citeproctext]{ref-thijssen20}
\CSLLeftMargin{{[}3{]} }%
\CSLRightInline{B. J. Thijssen, E. A. M. Klumperink, P. Quinlan, and B.
Nauta, {``30.4 a 370\(\mu\)w 5.5dB-NF BLE/BT5.0/IEEE 802.15.4-compliant
receiver with \textgreater63dB adjacent channel rejection at
\textgreater2 channels offset in 22nm FDSOI,''} in \emph{2020 IEEE
international solid-state circuits conference - (ISSCC)}, 2020, pp.
466--468, doi:
\href{https://doi.org/10.1109/ISSCC19947.2020.9062973}{10.1109/ISSCC19947.2020.9062973}.
}

\bibitem[\citeproctext]{ref-garvik19}
\CSLLeftMargin{{[}4{]} }%
\CSLRightInline{H. Garvik, C. Wulff, and T. Ytterdal, {``A 68 dB SNDR
compiled noise-shaping SAR ADC with on-chip CDAC calibration,''} in
\emph{2019 IEEE asian solid-state circuits conference (a-SSCC)}, 2019,
pp. 193--194, doi:
\href{https://doi.org/10.1109/A-SSCC47793.2019.9056925}{10.1109/A-SSCC47793.2019.9056925}.
}

\bibitem[\citeproctext]{ref-wulff17}
\CSLLeftMargin{{[}5{]} }%
\CSLRightInline{C. Wulff and T. Ytterdal, {``A compiled 9-bit 20-MS/s
3.5-fJ/conv.step SAR ADC in 28-nm FDSOI for bluetooth low energy
receivers,''} \emph{IEEE Journal of Solid-State Circuits}, vol. 52, no.
7, pp. 1915--1926, 2017, doi:
\href{https://doi.org/10.1109/JSSC.2017.2685463}{10.1109/JSSC.2017.2685463}.
}

\bibitem[\citeproctext]{ref-razavi20}
\CSLLeftMargin{{[}6{]} }%
\CSLRightInline{B. Razavi,
\emph{\href{https://doi.org/10.1017/9781108626200}{Design of CMOS
phase-locked loops}}. Cambridge University Press, 2020. }

\end{CSLReferences}
